\chapter{Tooling} \label{app:tooling}
%% DRAFT (Claude, 2026-10-13): condensed from README.md, main.py, config/config.yaml
%% and docker-compose.yml. Values (models, k, budgets) are deliberately left out:
%% they live in config.yaml and change too often to freeze here.
\TODO{REVIEW: Claude's draft of this appendix. Read it and rewrite it in your
own words.}

The code is public at \url{https://github.com/afonsoni/thesis}.

This appendix is a map, not a manual: enough to find your way around \sysname{}
and reproduce its numbers. The repository's README is the source of truth.

\section{Repository layout} \label{sec:tooling-layout}

\begin{table}[h]
    \centering
    \small
    \begin{tabularx}{\textwidth}{|l|X|}
        \hline
        \textbf{Path} & \textbf{Holds} \\
        \hline
        \texttt{main.py} & The command-line entry point (Section~\ref{sec:tooling-cli}). \\
        \texttt{config/config.yaml} & Every setting, in one file (Section~\ref{sec:tooling-config}). \\
        \texttt{src/rag/indexing/} & Parse, clean, chunk, classify, contextualise, embed, store. \\
        \texttt{src/rag/retrieval/} & Query analysis, expansion, routing, hybrid search, reranking. \\
        \texttt{src/rag/generation/} & Prompts and answer synthesis. \\
        \texttt{src/rag/evaluation/} & Variant ladder, metrics, LLM judge, leaderboard. \\
        \texttt{src/rag/common/} & Shared schemas, party and topic tables, the OpenRouter client. \\
        \texttt{data/} & Manifesto PDFs, their markdown, classified chunks, the gold set and results. \\
        \texttt{scripts/} & Gold-set anchoring, judge validation, maintenance. \\
        \texttt{tests/} & Unit, integration and end-to-end tests. \\
        \hline
    \end{tabularx}
    \caption{Top-level layout of the \sysname{} repository.}
    \label{tab:tooling-layout}
\end{table}

\section{Setup} \label{sec:tooling-setup}

Two Docker containers run PostgreSQL~17 with pgvector: the main database on
port 5432 and a throwaway test database on 5433. Python is 3.13, with every
dependency pinned in \texttt{requirements.txt}. The only secret is an
OpenRouter key, read from a \texttt{.env} file.

{\small
\begin{verbatim}
docker compose up -d
python3 -m venv .venv
.venv/bin/python -m pip install -r requirements.txt
echo "OPENROUTER_API_KEY=..." > .env
\end{verbatim}
}

Run everything through \texttt{.venv}. macOS ships its own \texttt{python3}
without any of these packages, and it fails at the first import.

\section{Commands} \label{sec:tooling-cli}

\begin{table}[h]
    \centering
    \small
    \begin{tabularx}{\textwidth}{|l|X|}
        \hline
        \textbf{Command} & \textbf{Does} \\
        \hline
        \texttt{parse} & Converts the manifesto PDFs to markdown with \texttt{marker}. \\
        \texttt{clean} & Rebuilds the header structure of the raw markdown with an LLM, rejecting any output that adds a word the source lacks. \\
        \texttt{index} & Chunks, classifies, contextualises, embeds and stores the corpus. Can stop after classification and resume from that checkpoint. \\
        \texttt{query} & Answers one question, with numbered citations. \\
        \texttt{eval} & Runs variants of the ladder over the gold set and refreshes the leaderboard. \\
        \hline
    \end{tabularx}
    \caption{Subcommands of \texttt{main.py}.}
    \label{tab:tooling-cli}
\end{table}

The cleaned markdown and the classified chunks are committed, so reproducing
the results starts at indexing; \texttt{parse} and \texttt{clean} are only
needed to rebuild the corpus from the PDFs. Four manifestos were cleaned by
hand and four by \texttt{clean}. \texttt{index --all} classifies every chunk
again with the LLM; \texttt{--from-classified} loads a party's committed
checkpoint instead and skips those calls. From the committed corpus to a first
answer, and from there to a headline number:

{\small
\begin{verbatim}
python main.py index --all
python main.py query "Quais são as propostas do PS para a educação?"

python scripts/anchor_gold.py candidates   # then review ANCHORING.md
python scripts/anchor_gold.py compile
python main.py eval --split test --curated-only
\end{verbatim}
}

\section{Database schema} \label{sec:tooling-schema}
%% Moved from the Proposition's vector store subsection (2026-10-25).

Each chunk is one row holding its text, the parent
section text, party and ideology, topic labels as an array, the heading path
as JSON, and a 1{,}024-dimensional vector. Similarity search uses an HNSW
index~\citep{malkov2020hnsw} on cosine distance
($m = 16$, $\mathit{ef\_construction} = 64$). Topic filters use a GIN index
on the label array. A generated \texttt{tsvector} column, computed with
PostgreSQL's Portuguese text-search configuration and also GIN-indexed,
serves lexical search. Because the column is generated by the database, it
cannot drift out of sync with the text it indexes. PostgreSQL runs in Docker,
so the whole store is reproducible from one configuration file.

\section{Configuration} \label{sec:tooling-config}

\texttt{config/config.yaml} has seven sections: \texttt{paths},
\texttt{embeddings}, \texttt{postgresql}, \texttt{llm} (models, fallbacks,
temperature, output and prompt budgets), \texttt{retrieval} (expansion count,
RRF constant, $k$, reranker), \texttt{indexing} (chunk sizes, classification and
context models) and \texttt{evaluation} (gold set, judge, default variants).
The loader is strict: a key the code reads but the file lacks raises an error,
instead of falling back to a default buried in the code.

\section{Tests} \label{sec:tooling-tests}

{\small
\begin{verbatim}
docker compose up -d postgres_test
python -m pytest              # unit + integration (test DB only)
python -m pytest tests/e2e    # live OpenRouter; spends quota
\end{verbatim}
}

The first two levels run before every commit. The end-to-end tests call the
real models, so they run only on request.
